% =====================================================================
% EveNet feasibility note for low-statistics NSBI in HH->bbtautau/bbgammagamma
% Working-note conventions: British -ise spellings, first-occurrence
% abbreviation expansion, \toverify{} flags for unverified content.
% Build: pdflatex twice (TOC + refs). No shell-escape, no external assets.
% =====================================================================
\documentclass[11pt]{article}

\usepackage[margin=2.5cm]{geometry}
\usepackage{amsmath,amssymb}
\usepackage{booktabs}
\usepackage{longtable}
\usepackage{enumitem}
\usepackage{xcolor}
\usepackage{microtype}
\emergencystretch=1.5em
\usepackage[colorlinks=true,linkcolor=blue!60!black,citecolor=blue!60!black,urlcolor=blue!60!black]{hyperref}

% ---- macros (match fm_nsbi_note conventions) ----
\newcommand{\kl}{\ensuremath{\kappa_\lambda}}
\newcommand{\bbtt}{\ensuremath{b\bar b\,\tau^+\tau^-}}
\newcommand{\bbgg}{\ensuremath{b\bar b\,\gamma\gamma}}
\newcommand{\tauh}{\ensuremath{\tau_h}}
\newcommand{\Nft}{\ensuremath{N_{\mathrm{ft}}}}
\newcommand{\mhh}{\ensuremath{m_{HH}}}
\newcommand{\mtt}{\ensuremath{m_{\tau\tau}}}
\newcommand{\mgg}{\ensuremath{m_{\gamma\gamma}}}
\newcommand{\toverify}[1]{\textcolor{red}{[verify: #1]}}

\title{Feasibility of EveNet for Low-Statistics NSBI in
$HH\to\bbtt$ and $HH\to\bbgg$:\\
Coverage Analysis, Tests, and Decision Plan}
\author{Working note --- \bbtt\ \kl\ NSBI project}
\date{10 July 2026 \\[2pt] \small v1 --- synthesised from arXiv:2601.17126, arXiv:2606.14870,
and the local EveNet integration harness (branch \texttt{fm})}

\begin{document}
\maketitle

\begin{abstract}
\noindent
We assess whether the EveNet event-level foundation model (FM)~\cite{Hsu:2026evenet}
can be fine-tuned to train high-precision, compute-efficient density-ratio
classifiers for a \kl\ measurement via neural simulation-based inference (NSBI)
in $HH\to\bbtt$ and $HH\to\bbgg$, with emphasis on the low-statistics regimes:
sparse signal Monte Carlo (MC) per coupling point (P1) and sparse
systematic-variation MC (P2). The central finding: \emph{neither hadronic taus
nor prompt photons exist in EveNet's object vocabulary or pre-training corpus},
and no supported extension path is documented. However, a three-tier analysis
(vocabulary / corpus / transferable structure) combined with an
information-flow analysis of the NSBI ratio decomposition shows the gap is
concentrated in identification information that approximately cancels in
exactly the low-statistics ratio components where FM data-efficiency is
needed. We therefore recommend proceeding with the already-built fine-tuning
harness (Option~A), upgraded with an ``unoccupied-corner'' type encoding and
calibration-closure metrics, before considering vocabulary surgery
(Option~B) or a new $\tau/\gamma$-aware FM (Option~C). A staged decision plan
with explicit gates and kill criteria is given, together with a ranked list of
open questions. The single largest scientific unknown --- whether fine-tuned FM
classifiers remain calibrated enough for density-ratio use --- is untested in
the literature and is answerable at near-zero marginal cost within our sweep.
\end{abstract}

\tableofcontents

% =====================================================================
\section{Scope and problem statement}
\label{sec:scope}

\paragraph{Task.} Measure the Higgs self-coupling modifier \kl\ via NSBI:
train classifiers whose outputs convert to per-event density ratios
$r(x;\kl,\kl')$ between process hypotheses, combined into an unbinned
likelihood. Signal discrimination between \kl\ hypotheses is driven by the
\mhh\ lineshape and the triangle--box interference (non-monotone in \kl,
with a near-degeneracy around $\kl\approx 2.45$), i.e.\ by event-level
kinematics rather than any single-object property.

\paragraph{Channels.} Primary: $HH\to\bbtt$ (2 $b$-jets, 2 taus
--- here focused on hadronic taus \tauh\ --- plus missing transverse
energy, MET, from neutrinos). Secondary increment: $HH\to\bbgg$
(2 $b$-jets, 2 prompt photons, no intrinsic MET from the photon leg).

\paragraph{The two low-statistics problems.}
\begin{description}[leftmargin=2.2em]
\item[P1] Sparse signal MC per coupling basis point:
  $\mathcal{O}(200\mathrm{k})$ events per $\kl\in\{0,1,5\}$, produced with
  Powheg $\to$ Pythia8 $\to$ Delphes $\Rightarrow$ high-variance ratio
  estimators.
\item[P2] Sparse systematic-variation MC $\Rightarrow$ difficulty
  parameterising $s(x;\kl,\theta)$ in nuisance parameters $\theta$;
  variation samples are typically $10\times$--$100\times$ smaller than
  nominal.
\end{description}

\paragraph{Quality gates already adopted by the analysis.}
Integral closure $<1\%$ per basis process and shape closure at the 5\% level
(\texttt{CLOSURE\_TOL}$\,=0.05$, treated as fixed --- a 5\% bias is already
significant). Any FM-based classifier must pass the \emph{same} gates as the
from-scratch baseline; area-under-curve (AUC)-type wins alone are
insufficient for likelihood-ratio use.

\paragraph{Compute framing.} NSBI multiplies training runs:
(processes) $\times$ (systematic variations) $\times$ (ensemble members)
$\times$ (\kl\ grid). Any FM must win on \emph{value per compute} across this
multiplicity, not on a single-classifier benchmark.

% =====================================================================
\section{Evidence base}
\label{sec:evidence}

\subsection{The Sophon null and the horizon/objective decomposition}
\label{sec:sophon}

Fine-tuning the single-jet tagger backbone Sophon
\toverify{Sophon arXiv ID; believed 2405.12972}~\cite{Li:2024sophon} for
$\kl=0$-vs-$5$ discrimination did \emph{not} beat from-scratch training
(scratch $\geq$ finetune $>$ frozen), and a cheap kinematic-ceiling
BDT/MLP matched all three. Diagnosis (per
\texttt{FOUNDATION\_MODEL\_BRIEF.md}): two independent mismatches,
\begin{enumerate}[label=(\roman*)]
\item \textbf{horizon} --- single-jet pretext vs.\ event-level physics
  (the \kl\ signal is a multi-object correlation), and
\item \textbf{objective} --- discrete flavour classification vs.\ a
  continuous, interference-driven ratio.
\end{enumerate}
EveNet is the cleanest available experiment for \emph{disentangling} these:
it is event-level but \emph{not} coupling-aware. An EveNet null would imply
the objective mismatch is binding; an EveNet win would implicate the horizon
mismatch as the cause of the Sophon null. This disentangling logic is the
primary scientific motivation for the EveNet test, independent of the
$\tau/\gamma$ question.

\subsection{Controlled pre-training-objective study (arXiv:2606.14870)}
\label{sec:elsharkawy}

Elsharkawy et al.~\cite{Elsharkawy:2026} compare pre-training objectives on a
fixed backbone (OmniLearned Point--Edge Transformer, PET; 100k/7M/51M
parameters) pre-trained on JetClass (100M jets; MadGraph5+Pythia8+Delphes
3.4.3)~\cite{Qu:2022part}, fine-tuned on top tagging with
$\Nft\in\{10^3,10^4,10^5,1.2\times10^6\}$ and on JetNet generation. Findings
(labels F1--F4 as in the companion FM/NSBI note):

\begin{description}[leftmargin=2.2em]
\item[F1] \textbf{Classifier+MPM wins the low-label regime.} At
  $\Nft=10^4$ the classifier + masked-particle-modelling (MPM) hybrid is the
  top configuration at every model size (by their Data-Efficiency Score);
  with plentiful labels ($\Nft=1.2$M) pure classifier pre-training is best.
  ``Low-label'' means few labelled downstream training events; in
  simulation-based HEP every event is labelled for free, so for us it simply
  means \emph{low-statistics fine-tuning samples}. Our P2 samples
  ($10^3$--$10^4$) sit squarely in the winning regime; our 200k/\kl-point
  signal samples sit near the boundary where the advantage starts narrowing.
\item[F2] \textbf{Generative pre-training does not transfer to
  classification.} Pure flow-matching pre-training matches or underperforms
  from-scratch at all sizes.
\item[F3] \textbf{Classification and generation are orthogonal.}
  Flow matching must be in the pre-training mix to obtain downstream
  generative advantage, and vice versa.
\item[F4] \textbf{Pre-training validation loss is an imperfect downstream
  proxy} (Pearson $r\approx-0.75$); downstream performance peaks
  \emph{before} pre-training converges. Checkpoint selection is an open
  problem --- for NSBI it becomes a \emph{calibration-selection} problem.
\end{description}

Two additional results matter directly for the present decision:
\begin{itemize}
\item \textbf{Pre-training data volume is nearly irrelevant at low labels.}
  In their $10^5\!\to\!10^8$ pre-training-set scan (Micro model), scaling
  slopes with scarce downstream labels are essentially flat
  ($\bar m = 1.4$, vs.\ $\bar m = 182.5$ with abundant labels); only
  standalone MPM retains a positive slope. \emph{Objective choice, not
  corpus size, drives the low-label gain.} This substantially reduces the
  corpus-production cost of any continued-pre-training option
  (Sec.~\ref{sec:options}).
\item \textbf{Two evidence gaps.} (a) \emph{Calibration is never tested}:
  all classification metrics are rejection-type ($1/\varepsilon_b$ at
  $\varepsilon_s=0.3$). Whether a fine-tuned FM classifier converts to an
  unbiased density ratio is unknown. (b) \emph{No genuinely out-of-domain
  transfer is tested}: JetClass$\to$top-tagging is class-overlapping (tops
  are in the pre-training corpus); no downstream task involves object types
  absent from pre-training. Both gaps sit exactly on our critical path.
\end{itemize}

\subsection{EveNet (arXiv:2601.17126): what it is}
\label{sec:evenet}

\subsubsection{Input schema}
Per-object token = 7 features:
$(E,\,p_T,\,\eta,\,\phi,\,\texttt{btag},\,\texttt{isLepton},\,\texttt{charge})$,
with $E,p_T$ log-normalised. Global (per-event) conditioning = 10 features:
$(\mathrm{MET},\,\mathrm{MET}_\phi,\,n_\ell,\,n_b,\,n_j,\,H_T,\,H_T^\ell,\,
M_{\mathrm{all}},\,M_{\ell\ell},\,M_{bb})$. \textbf{MET is not a token} ---
it enters only through the global block. There is \emph{no type-embedding
vocabulary}: object identity is carried entirely by the three discrete flags.
There is no \texttt{isTau}, no \texttt{isPhoton}, no \texttt{isMET}, and no
extensible type table. The schema is fixed jointly by the
\texttt{event\_info} YAML used in preprocessing \emph{and} training; changing
it changes the input dimension of the \texttt{GlobalEmbedding}/input blocks
and breaks strict checkpoint compatibility.

\subsubsection{Pre-training corpus}
$\sim$33 billion simulated events, $\sim$500M retained after preselection.
Processes (their Table~1): QCD multijet (2--4$j$); $t\bar t$
(hadronic/1L/2L $+0,1j$); $t\bar tW$, $t\bar tZ$; $W$+jets, $Z$+jets
($0,1,2j$); dibosons $WW$, $ZZ$, $WZ$; and SM $H\to WW^*$ (0L/1L/2L).
Chain: MadGraph5\_aMC@NLO $\to$ Pythia $\to$ Delphes with a \emph{generic}
detector card (explicitly neither CMS nor ATLAS). ``Leptons'' throughout
means electrons or muons.
\textbf{No prompt-photon process (no $H\to\gamma\gamma$, $\gamma$+jets,
$t\bar t\gamma$) and no hadronic-tau final state appears in the corpus.}
Hadronic taus appear in the paper only as an event \emph{veto} in the
$X\to YH$ selection, never as input tokens.

\subsubsection{Pre-training objectives}
Two-stage, diffusion-time-conditioned:
\begin{itemize}
\item \textbf{Stage 1 (SSL only):} masked-particle reconstruction via
  diffusion inpainting (mask a subset of tokens; denoise conditioned on the
  rest), $\mathcal{L}=\lVert v - v_{\mathrm{pred}}\rVert_2^2$, masking
  probability ramped during training.
\item \textbf{Stage 2 (multi-task):} classification (process discrimination,
  noise-weighted $\alpha^2(t)\,\mathrm{CE}$, decay-channel variations
  \emph{absorbed within} each class), SSL generation, and \emph{supervised
  generation = neutrino reconstruction} (denoise injected tokens for
  undetectable components conditioned on observed objects). Losses summed
  directly, no gradient surgery. Segmentation and assignment heads exist but
  are fine-tune-only.
\end{itemize}
\textbf{No systematics-aware or coupling/nuisance-parameterised component
exists in the pre-training objective.}

\subsubsection{Fine-tuning interface and demonstrated results}
Three initialisations are compared throughout: \emph{Scratch},
\emph{SSL} (Stage-1 weights), \emph{Full} (Stage-1+2). Standard recipe:
joint fine-tune with backbone learning rate $=10\%$ of head learning rate;
optional adapters between transformer blocks (trainable parameters
$\sim$20M$\to$8M). Demonstrated tasks:
\begin{enumerate}[label=(\alph*)]
\item \textbf{$X\to YH\to b\bar b\,WW^*\to b\bar b\,\ell\nu qq'$}
  (CMS Open Data full-sim, 121 mass points, 1k--10k signal events each):
  Full $>$ Scratch by $\sim$50\% in significance-improvement characteristic
  (SIC); $>$ XGBoost by $\sim$20\%, $>$ TabPFN by $\sim$10\%;
  $\sim3\times$ faster convergence.
\item \textbf{Exotic $H\to aa\to 4b$} ($m_a=30$\,GeV; explicitly absent from
  the corpus): SIC 4.1 (Full) vs.\ 1.6 (Scratch) vs.\ 1.4 (SPANet); at
  \emph{5\% of training data} SIC 1.7, exceeding both baselines at full
  data; pairing efficiency 58\% vs.\ 26\%.
\item \textbf{$t\bar t$ dilepton spin correlations} (scan 1.5\%--150\% of
  2.83M events): at 1.5\%, $\Delta D = 1.85\%$ (Full) vs.\ 2.68\% (Scratch);
  lepton--quark pairing 48\% vs.\ 24\%; extraction reported stable under
  JES-type shifts.
\item \textbf{$\Upsilon\to\mu\mu$ rediscovery in real CMS Open Data} (weak
  supervision, $\sim$12k events): median significance $7.6\pm0.4\sigma$
  (Full) vs.\ $1.45\sigma$ (Scratch) vs.\ $6.4\sigma$ published benchmark
  --- the starkest low-statistics pre-training win.
\end{enumerate}
Cautionary result: \emph{mass-parameterised} fine-tuning showed
``systematically reduced sensitivity'' (their explanation: kinematic
diversity across the mass grid limits parameter sharing). This warns
against a \kl-parameterised head as the first NSBI design; our
per-basis-point + linear-combination construction sidesteps it.

\subsubsection{Architecture, compute, availability}
PET encoder, 18.8M parameters (hidden 256, 16 layers, 16 heads, 10 local
$k$NN layers); task heads 1.3--8.7M. Pre-trained on 512 A100s (Perlmutter),
LION + EMA, PyTorch~2.6 + Lightning + Ray. Code is public at
\texttt{github.com/UW-EPE-ML/EveNet\_Public}; checkpoints are public at
\texttt{huggingface.co/Avencast/EveNet}
(\texttt{checkpoints.20M.a4.last.ckpt} and \texttt{SSL.20M.last.ckpt});
a \texttt{network-100M.yaml} scaled template
exists in the repository. Scope restriction: resolved regime only (no
boosted/large-$R$ topologies). Robustness to the generic-Delphes$\to$CMS
full-sim domain shift is demonstrated \emph{empirically}, not by
construction.

% =====================================================================
\section{The $\tau/\gamma$ coverage gap: three-tier analysis}
\label{sec:threetier}

``Does EveNet cover $\tau/\gamma$ final states?'' decomposes into three
separable questions with different remediation costs:

\begin{table}[h]
\centering\small
\begin{tabular}{@{}p{3.2cm}p{5.6cm}p{2.2cm}p{3.4cm}@{}}
\toprule
\textbf{Tier} & \textbf{Finding} & \textbf{Verdict} & \textbf{Remedy if binding} \\
\midrule
Vocabulary (can the schema represent \tauh/$\gamma$?) &
No type table exists; identity = 3 flags
(\texttt{btag}, \texttt{isLepton}, \texttt{charge}).
No $\tau/\gamma$ flag. MET never a token. &
Gap real, but shallower than feared (Sec.~\ref{sec:encoding}) &
Corner encoding (free) or input-widening surgery (cheap--medium) \\
\addlinespace
Corpus (did pre-training see $\tau/\gamma$ physics?) &
Zero prompt-photon processes; zero \tauh\ final states;
\tauh\ used only as an event veto. &
Gap real &
Continued pre-training on $\tau/\gamma$-enriched corpus (medium; volume
requirement reduced by the flat low-label scaling result of
\cite{Elsharkawy:2026}) \\
\addlinespace
Transferable structure (does the learned representation help anyway?) &
All demonstrated out-of-corpus successes ($H\to aa\to4b$, $X\to YH$,
fast-sim$\to$full-sim, sim$\to$data) are \emph{new topologies of known
object types}. The truly out-of-vocabulary case is untested anywhere in
the literature. &
\textbf{Open --- the decisive question} &
Not remediable by construction; must be \emph{measured} (Gate G0,
Sec.~\ref{sec:plan}) \\
\bottomrule
\end{tabular}
\caption{Three-tier decomposition of the $\tau/\gamma$ coverage gap.}
\label{tab:tiers}
\end{table}

\noindent
The trap to avoid is collapsing these into a binary ``EveNet doesn't do taus
$\Rightarrow$ build our own FM.'' Tiers 1--2 are established facts with cheap
or medium remedies; tier 3 is an empirical question our existing harness can
answer directly.

% =====================================================================
\section{Why the gap is likely non-binding: ratio-component information asymmetry}
\label{sec:asymmetry}

The NSBI likelihood decomposes into per-process density ratios, and the
$\tau/\gamma$ \emph{identification} information is asymmetrically distributed
across them:

\begin{description}[leftmargin=2.2em]
\item[Signal--signal ratios $r(x;\kl,\kl')$.] Numerator and denominator
  share an \emph{identical final state}. Type-identification quantities
  (DeepTau-like scores, decay modes, photon ID) are (to good approximation)
  common factors that cancel in the ratio. The \kl\ dependence lives in the
  \mhh\ lineshape and triangle--box interference --- functions of the object
  four-vectors, which the stock token $(E,p_T,\eta,\phi)$ carries
  \emph{exactly}.
\item[Signal-vs-background discriminators.] Here $\tau$-ID/$\gamma$-ID
  genuinely matter (fake rejection against $t\bar t$, QCD, $\gamma$+jets).
  But these components are \emph{MC-rich} --- $t\bar t$ is the most abundant
  sample in the analysis --- so they are not where FM data-efficiency is
  needed, and existing full-feature classifiers can be retained for them.
\item[Systematic-variation ratios (nominal/varied, P2).] Variation samples
  share the nominal final state; type-ID again approximately cancels, and
  energy-scale-type shape effects are kinematic and token-visible.
\end{description}

\noindent\textbf{Consequence.} The ratio components where the FM's low-label
advantage is needed (P1 signal ratios, P2 variation ratios) are precisely the
components where the missing $\tau/\gamma$ vocabulary costs the least. The
stock-schema workaround already encoded in the harness (\tauh\ $\to$ generic
jet) is approximately lossless \emph{for the specific low-statistics question
the FM is being asked to answer}, with the losses concentrated in MC-rich
background rejection.

\paragraph{Honest caveats.}
\begin{enumerate}[label=(\roman*)]
\item ``Fully preserved'' applies to \emph{visible} kinematics. In \bbtt,
  tau neutrinos smear visible \mhh; SVfit/FastMTT \mtt\ --- which partially
  recovers the neutrino momenta --- has no slot in the stock schema and
  \emph{is} genuinely lost. Mitigations: MET (the aggregate neutrino
  footprint) is in the global conditioning, and EveNet's Stage-2 pre-training
  explicitly included \emph{neutrino reconstruction as a supervised denoising
  task} --- the backbone was trained to infer invisible components from
  visible tokens + MET. Whether this recovers what FastMTT provides is an
  empirical question the sweep answers (open question Q5).
\item The cancellation argument is approximate: selection efficiencies and
  fake contamination differ slightly between \kl\ hypotheses through
  kinematic correlations with ID variables. Second-order, but the closure
  gates will price it.
\item For \bbgg\ the preservation argument is \emph{stronger} than for
  \bbtt: photons produce no neutrinos, so visible kinematics = full
  kinematics and \mgg\ is exactly computable from the tokens. Conversely the
  FM \emph{headroom} is smaller (the \mhh\ observable is already cleaner),
  which is why \bbgg\ is sequenced as an increment, not the primary test.
\end{enumerate}

% =====================================================================
\section{Encoding strategies within the fixed schema}
\label{sec:encoding}

\subsection{Status quo: anonymous-jet encoding}
The current adapter (\texttt{scripts/delphes\_to\_evenet\_npz.py}) emits
\tauh\ as a generic jet: (\texttt{btag}, \texttt{isLepton},
\texttt{charge}) $= (0,0,0)$. Kinematics exact; $\tau$ identity, charge, and decay mode
discarded; the network must infer $\tau$-ness from kinematics alone.
Documented, deliberate, checkpoint-safe.

\subsection{Unoccupied-corner encoding (proposed, zero surgery)}
The three flags span a $2\times2\times3=12$-point discrete lattice; every
token's \emph{type} is a lattice corner. In EveNet's pre-training corpus only
four corners ever occur:

\begin{table}[h]
\centering\small
\begin{tabular}{@{}llll@{}}
\toprule
\textbf{Corner} (\texttt{btag}, \texttt{isLep}, \texttt{chg}) &
\textbf{Corpus occupant} & \textbf{Proposed reassignment} & \textbf{Status} \\
\midrule
$(0,0,0)$   & light jet        & --- & occupied \\
$(1,0,0)$   & $b$-jet          & --- & occupied \\
$(0,1,\pm1)$& $e^\pm/\mu^\pm$  & --- & occupied \\
$(0,0,\pm1)$& \emph{never seen} & \textbf{\tauh} (``charged jet''), with
  measured charge & vacant \\
$(0,1,0)$   & \emph{never seen} & \textbf{$\gamma$} (``neutral lepton'') & vacant \\
$(1,1,\cdot)$, $(1,0,\pm1)$ & \emph{never seen} & reserved & vacant \\
\bottomrule
\end{tabular}
\caption{The flag lattice: occupied vs.\ vacant corners, and the proposed
$\tau/\gamma$ assignments. Values are in-range for the pre-trained
normalisers; the checkpoint loads unchanged.}
\label{tab:corners}
\end{table}

\noindent
This gives the network a \emph{distinguishable} type signal for \tauh\
(recovering its measured charge, currently discarded) and $\gamma$ inside the
existing 7 dimensions --- no weight surgery, full checkpoint compatibility,
a one-line adapter change per object type.

\subsection{Risks of the corner encoding (candidate failure mechanisms)}
\label{sec:cornerrisks}
The backbone has no prior \emph{about} vacant corners; ``strictly more
information at the data level'' does not guarantee a better classifier after
a \emph{finite} fine-tune. Candidate mechanisms by which it could hurt,
all measurable by an A/B test (Gate G1):
\begin{enumerate}[label=(R\arabic*)]
\item \textbf{Out-of-distribution embedding:} inputs at never-seen flag
  combinations land the token embedding in unexercised regions of
  representation space; early fine-tuning could be slower or less stable
  (mitigated by the backbone-LR$=10\%$ recipe).
\item \textbf{Shortcut/overfit capacity at very low \Nft:} an extra clean
  discrete bit gives the network memorisation shortcuts; at
  $\Nft\sim10^3$ this could \emph{worsen} generalisation relative to the
  anonymous encoding.
\item \textbf{Attention misrouting:} pre-trained heuristics keyed on
  $\texttt{charge}\neq0\Rightarrow$ lepton may initially misroute \tauh\
  tokens through lepton-like attention patterns (transient, but could
  matter in short fine-tunes).
\end{enumerate}

\subsection{What cannot be encoded without surgery}
$\tau$-ID score (DeepTau-like), decay mode, photon-ID quality, and
FastMTT/SVfit \mtt\ have no slot. Adding any of them widens the token (or the
global block) and breaks direct checkpoint loading --- that is Option~B
territory (Sec.~\ref{sec:options}). Overwriting one of the 10 global
conditions with \mtt\ is technically possible but abuses a pre-trained
normaliser and is not recommended.

% =====================================================================
\section{Current integration state (branch \texttt{fm})}
\label{sec:harness}

Already built and locally verified, in
\texttt{examples/HH\_bbtautau\_kappalambda\_sophon/}:

\begin{itemize}
\item \textbf{Adapter} \texttt{scripts/delphes\_to\_evenet\_npz.py}:
  Delphes ROOT $\to$ NPZ with \texttt{x} of shape $(N,16,7)$ (zero-padded,
  $p_T$-sorted), \texttt{conditions} $(N,10)$, sequential-count mask,
  event weights (NLO signs preserved), binary class label. Non-negativity
  clipping matches EveNet's \texttt{log1p} scaling. \texttt{--smoke} mode
  fabricates schema-valid events without \texttt{uproot}.
\item \textbf{Systematic injector} \texttt{scripts/inject\_systematic.py}:
  analytic distortions of jet tokens --- \texttt{jes} (constant scale
  $\times(1+\alpha)$) and \texttt{jes\_mhh}
  ($\alpha\tanh((M_{\mathrm{all}}-m_0)/w)$, a shape systematic correlated
  with the \kl\ observable) --- with globals recomputed. Closure is exact by
  construction, so P2 value is \emph{measurable} without real variation MC.
  Known TODO: JES$\to$MET propagation deliberately omitted in v1.
\item \textbf{Sweep generator} \texttt{scripts/make\_klambda\_configs.py}:
  3 modes (finetune / frozen / scratch) $\times$ 5 dataset fractions
  $(0.01,0.03,0.1,0.3,1.0)$ $\times$ 5 seeds $=75$ configurations +
  launcher script; adapted from the EveNet Exotic-Higgs study template.
\item {\sloppy \textbf{Configs}: \texttt{workflow\_klambda.yaml} (control
  file; paths are placeholders) and \texttt{event\_info\_klambda.yaml},
  whose INPUTS block is byte-identical to EveNet's \texttt{pretrain.yaml}
  so the 20M backbone binds (assignment/permutation heads present but
  off).\par}
\item \textbf{Evaluation} \texttt{scripts/plot\_data\_efficiency.py}:
  weighted AUC vs.\ training fraction, mean$\pm$std over seeds, with
  kinematic-ceiling overlay. \emph{Currently AUC-only} --- the closure-metric
  upgrade (Gate G0.5) is the key missing piece.
\item \textbf{Runtime}: Perlmutter shifter image
  \texttt{avencast1994/evenet:1.5}; checkpoints via
  \texttt{hf download Avencast/EveNet}.
\end{itemize}

\noindent Engineering TODOs carried from the harness handoff: create
\texttt{options\_frozen.yaml}; verify the prediction key
\texttt{classification/klambda}; load-test the configs inside the shifter
image before launching the 75 jobs; and verify Delphes branch names
against real files.

% =====================================================================
\section{Options and costs}
\label{sec:options}

\begin{table}[h]
\centering\small
\begin{tabular}{@{}p{2.6cm}p{3.8cm}p{4.2cm}p{3.6cm}@{}}
\toprule
& \textbf{A. Fine-tune stock EveNet} (+ corner encoding) &
\textbf{B. Vocabulary surgery + continued pre-training} &
\textbf{C. New $\tau/\gamma$-aware FM} \\
\midrule
What it is &
Existing 75-config harness; optional corner encoding
(Sec.~\ref{sec:encoding}) &
Widen token to 8--9 dims (\texttt{isTau}/\texttt{isPhoton} or $\tau$-ID
score); zero-init new input columns (function-preserving at init);
continue SSL+classification pre-training on a $\tau/\gamma$-enriched
Delphes corpus ($Z\to\tau\tau$, $H\to\tau\tau$, $\gamma$+jets,
$H\to\gamma\gamma$, $t\bar t\gamma$, HH grid) &
Full pre-training from scratch with $\tau/\gamma$ in vocabulary and
corpus \\
\addlinespace
Weight reuse & Full & Full except first linear layer & None \\
\addlinespace
New corpus & None &
Modest: F-volume result of \cite{Elsharkawy:2026} (flat low-label
scaling) suggests $\mathcal{O}(10^7)$ events from the existing Delphes
pipeline plausibly suffices &
EveNet-scale was 33B raw / 500M retained;
$\mathcal{O}(10^2$--$10^3)\times$ A--B \\
\addlinespace
Compute & Fine-tune only: hours--days on Perlmutter &
+ continued pre-training of an 18.8M model: days--weeks &
Months + generation campaign (EveNet used 512 A100s) \\
\addlinespace
Key risk & $\tau/\gamma$ information ceiling (bounded by
Sec.~\ref{sec:asymmetry}) &
No established HEP recipe for continued pre-training; representation
drift; schema fork diverges from upstream &
Duplicates $\sim$1 year of engineering; \textbf{and if A nulls, C likely
nulls too} (Sec.~\ref{sec:plan}) \\
\bottomrule
\end{tabular}
\caption{The three strategic options.}
\label{tab:options}
\end{table}

\noindent
On Option C, note also the existence of TauFM~\cite{TauFM:2025} (jet-FM-based
$\tau$ reconstruction) as adjacent prior art, and that the EveNet group
invites process contributions and ships a 100M-parameter template ---
if Option C is ever reached, upstream collaboration dominates a solo rebuild.

% =====================================================================
\section{Staged decision plan}
\label{sec:plan}

Cheapest-information-first. Each gate has an explicit pass condition and a
kill/branch interpretation.

\begin{description}[leftmargin=2.6em,style=nextline]
\item[G0 --- Run the existing sweep (Option A, stock or corner schema).]
  \emph{Pass:} finetune $>$ scratch beyond the compute penalty
  \textbf{and} beats the kinematic-ceiling baseline on value-per-compute.
  \emph{Interpretation of a null:} event-level pre-training did not fix the
  Sophon null $\Rightarrow$ the \emph{objective} mismatch is binding
  (Sec.~\ref{sec:sophon}); no amount of $\tau$-vocabulary repairs an
  objective mismatch, so \textbf{a null here also condemns Option C in its
  like-for-like form} --- the missing $\tau$-ID is not what limits
  \kl-separation (Sec.~\ref{sec:asymmetry}); the surviving FM directions
  would be coupling-aware pretexts, not coverage extensions.
\item[G0.5 --- Add calibration/closure metrics to the sweep evaluation.]
  Integral closure $<1\%$ per basis process and shape closure at
  \texttt{CLOSURE\_TOL}$=0.05$, alongside AUC, for every (mode, fraction,
  seed) cell. Near-zero marginal cost; fills the literature's largest gap
  (Sec.~\ref{sec:elsharkawy}); publishable in either direction; directly
  feeds the low-statistics-systematics ``money plot'' goal.
  \emph{Branch:} if FM fine-tunes win AUC but fail closure where scratch
  passes, the FM inductive bias is biasing the ratio --- a new, reportable
  failure mode, and a hard blocker for NSBI adoption regardless of G0.
  \emph{Decision primacy (fixed 2026-07-12, before any sweep curves existed):}
  the adoption metric is the closure \textbf{left-shift} --- finetune passing
  both gates at a smaller training fraction than scratch (e.g.\ 3\% vs.\ 30\%
  $=$ a $10\times$ equivalent-data multiplier) \emph{is} an adoption-grade win
  even at AUC parity; an AUC-only win with failing closure is not. AUC is a
  secondary guard (finetune must not materially regress below scratch).
  Calibration, not ordering, is the binding axis: the full-statistics ceiling
  measurement shows AUC saturating by the 10\% fraction (0.792 at 1\% vs.\
  0.817 at 100\%) while its closure gates first pass only at 10\%, with a
  method-limited $|\mathrm{IC}|$ floor of ${\approx}1.7\%$.
\item[G1 --- Corner-encoding A/B test (\tauh: anonymous vs.\ $(0,0,\pm1)$).]
  One adapter flag; run at two dataset fractions $\times$ 5 seeds.
  \emph{Pass:} corner $\geq$ anonymous at all fractions (watch R2 of
  Sec.~\ref{sec:cornerrisks} at the lowest fraction).
\item[G2 --- \bbgg\ increment.] Photons via $(0,1,0)$; only the adapter's
  object-filling changes. Tests whether the corner trick generalises to a
  second unseen type, on the channel where the kinematic-preservation
  argument is strongest.
\item[G3 --- Option B surgery (only if corners measurably saturate).]
  Zero-init input widening + continued SSL pre-training; corpus sized
  modestly per the flat-volume result. \emph{Pass:} beats the corner
  encoding by more than the added compute.
\item[G4 --- Option C (only if transfer is proven valuable \emph{and}
  surgery demonstrably fails).] Prefer collaboration with the EveNet
  group over a solo build.
\end{description}

\noindent\textbf{Standing constraint:} no $\sigma_{\kl}$ significance claims
from any branch until the upstream NSBI MC-statistics estimator critique
\cite{Shyamsundar:2025} is settled (the wifi-ensemble programme
\cite{Wifi:2025} is the candidate resolution and is deferred to a separate
project).

% =====================================================================
\section{Open questions, ranked}
\label{sec:openq}

Priority tags [P0/P1/P2] and effort tags [S/M/L] follow the note-family
convention.

\begin{enumerate}[label=\textbf{Q\arabic*}.]
\item \textbf{[P0, S] Calibration of fine-tuned FM classifiers.} Does a
  pre-trained backbone help, hurt, or not affect ratio closure at fixed
  AUC? Untested in \cite{Elsharkawy:2026} (rejection metrics only) and in
  \cite{Hsu:2026evenet} (SIC only). Answered by G0.5.
\item \textbf{[P0, S] Does event-level pre-training beat scratch at all
  for a continuous interference-driven ratio?} The
  horizon-vs-objective disentangling test (G0). All published EveNet wins
  are discrete classification/anomaly tasks.
\item \textbf{[P0, --- (blocking for claims)] MC-statistics estimator
  critique} \cite{Shyamsundar:2025}: settle before quoting any
  \kl\ significance from FM or non-FM pipelines.
\item \textbf{[P1, S] Corner vs.\ anonymous encoding} (G1), including the
  low-\Nft\ overfit risk R2.
\item \textbf{[P1, S] P2 value measurement.} Data-efficiency of
  nominal-vs-varied ratio learning under the analytic injector: does the FM
  advantage \emph{grow} as variation-sample size shrinks (the money plot),
  and does it survive the \texttt{jes\_mhh} shape systematic that is
  correlated with the \kl\ observable?
\item \textbf{[P1, M] Quantify the FastMTT/\mtt\ and $\tau$-ID loss.}
  Compare the kinematic ceiling with and without \mtt/ID features to bound
  what the stock schema forfeits in absolute terms; separately, test
  whether EveNet's neutrino-reconstruction pretext implicitly recovers
  \mtt-like information (probe the fine-tuned representation).
\item \textbf{[P1, S] Parameterised-head caution.} EveNet's
  mass-parameterised fine-tuning \emph{reduced} sensitivity; verify our
  per-basis-point + linear-combination NSBI construction is unaffected, and
  treat any \kl-parameterised head proposal as requiring its own gate.
\item \textbf{[P2, M] Continued-pre-training recipe stability} (for G3):
  zero-init widening, replay mixture of original-corpus processes vs.\
  $\tau/\gamma$ additions, representation-drift monitoring. No established
  HEP precedent.
\item \textbf{[P2, S] Domain gap: generic Delphes card vs.\
  \texttt{cms\_card\_v1}.} Empirically bridged in the paper
  (generic$\to$CMS full-sim); verify on our samples, since our card is
  Run-3-tuned (UParT/DeepTau baselines).
\item \textbf{[P2, M] Value-per-compute accounting across the NSBI
  multiplicity} (processes $\times$ systematics $\times$ ensembles
  $\times$ \kl\ grid), including whether a \emph{single} fine-tuned
  backbone can be shared across ratio components (amortisation would
  strongly favour the FM).
\end{enumerate}

% =====================================================================
\section{Summary verdict}
\label{sec:verdict}

\begin{itemize}
\item \textbf{Coverage:} EveNet does \emph{not} cover $\tau$ or $\gamma$
  final states --- no vocabulary slot, no corpus processes, no supported
  extension path. This is established fact, not conjecture
  (Sec.~\ref{sec:evenet}).
\item \textbf{Feasibility anyway:} the gap is concentrated in
  identification information that approximately cancels in the
  low-statistics ratio components (signal--signal and nominal--varied)
  where FM data-efficiency is actually needed; the \kl-carrying
  kinematics survive the stock schema exactly, up to the visible-kinematics
  caveat (\mtt\ loss, partially mitigated by the neutrino-reconstruction
  pretext) (Sec.~\ref{sec:asymmetry}).
\item \textbf{Do not build a new FM now, and do not start with surgery.}
  The existing harness is the correct next instrument: it runs the
  disentangling experiment (G0) whose null would condemn a like-for-like
  Option C anyway, and --- upgraded with closure metrics (G0.5) --- it
  resolves the literature's largest untested assumption for FM-based NSBI
  at near-zero marginal cost.
\item \textbf{Cheap upgrades first:} the unoccupied-corner encoding
  (\tauh\ as $(0,0,\pm1)$, $\gamma$ as $(0,1,0)$) recovers a type signal
  with zero weight surgery and one adapter line; it strictly dominates the
  anonymous-jet encoding at the data level, with finite-sample risks that
  G1 measures directly.
\item \textbf{Escalation is evidence-gated:} corners $\to$ zero-init
  widening + modest continued pre-training (volume de-risked by the flat
  low-label scaling result) $\to$ new/collaborative FM, each step only on
  demonstrated saturation of the previous one.
\end{itemize}

% =====================================================================
\begin{thebibliography}{99}\small

\bibitem{Hsu:2026evenet}
S.-C. Hsu et al.,
``EveNet: A Foundation Model for Particle Collision Data Analysis,''
arXiv:2601.17126 [hep-ex] (2026).
Code: \texttt{github.com/UW-EPE-ML/EveNet\_Public};
checkpoints: \texttt{huggingface.co/Avencast/EveNet}.

\bibitem{Elsharkawy:2026}
Y. Elsharkawy, J. Birk, V. Mikuni, W. Bhimji, G. Kasieczka, B. Nachman,
``Pre-Training for Simulation-Based Science: A Study on Jet Foundation
Model Training Objectives,''
arXiv:2606.14870 (2026).

\bibitem{Qu:2022part}
H. Qu, C. Li, S. Qian,
``Particle Transformer for Jet Tagging'' (ParT; introduces the JetClass
dataset),
arXiv:2202.03772 (2022).

\bibitem{Li:2024sophon}
C. Li et al.,
``Accelerating Resonance Searches via Signature-Oriented Pre-training''
(Sophon), \toverify{arXiv ID; believed 2405.12972} (2024).

\bibitem{Harris:2024rs3l}
P. Harris et al.,
``Re-Simulation-based Self-Supervised Learning for Pre-Training Foundation
Models'' (RS3L),
arXiv:2403.07066; Phys.\ Rev.\ D 111 (2025) 032010.

\bibitem{Golling:2024mpm}
T. Golling et al.,
``Masked Particle Modeling on Sets: Towards Self-Supervised High Energy
Physics Foundation Models'' (MPM),
arXiv:2401.13537 (2024).

\bibitem{Birk:2024omnijet}
J. Birk, A. Hallin, G. Kasieczka,
``OmniJet-$\alpha$: The first cross-task foundation model for particle
physics,''
arXiv:2403.05618; Mach.\ Learn.: Sci.\ Technol.\ 5 (2024) 035031.

\bibitem{Mikuni:2024omnilearn}
V. Mikuni, B. Nachman,
``OmniLearn: A Method to Simultaneously Facilitate All Jet Physics Tasks,''
arXiv:2404.16091 (2024).

\bibitem{Dillon:2021jetclr}
B. M. Dillon et al.,
``Symmetries, Safety, and Self-Supervision'' (JetCLR),
arXiv:2108.04253; SciPost Phys.\ 12 (2022) 188.

\bibitem{HEPJEPA:2025}
``HEP-JEPA: A foundation model for collider physics using joint embedding
predictive architecture,''
arXiv:2502.03933 (2025).

\bibitem{TauFM:2025}
``Tau reconstruction with a jet foundation model,''
arXiv:2503.19165 (2025).

\bibitem{Shyamsundar:2025}
P. Shyamsundar et al.,
critique of MC-statistical uncertainty estimation in NSBI,
arXiv:2505.19156 (2025).

\bibitem{Wifi:2025}
``Wifi ensembles'' MC-statistics resolution programme,
arXiv:2506.00113 (2025).

\end{thebibliography}

\end{document}
